% Espanol (Mexico) con XeLaTeX: fontspec para fuentes Unicode, polyglossia
% para la division silabica y los nombres del documento en la variante mexicana.
\usepackage{fontspec}
\usepackage{polyglossia}
\setdefaultlanguage[variant=mexican]{spanish}
% polyglossia no traduce los nombres de listings.
\AtBeginDocument{\providecommand{\lstlistingname}{}\renewcommand{\lstlistingname}{Listado}%
  \providecommand{\lstlistlistingname}{}\renewcommand{\lstlistlistingname}{Índice de listados}}
\usepackage{amsmath,amssymb}
\usepackage{graphicx}
\usepackage[margin=2.35cm]{geometry}
\usepackage[most]{tcolorbox}
\usepackage{etoolbox}
\usepackage{listings}
\usepackage{hyperref}
\usepackage{booktabs}
\usepackage{array}
\usepackage{float}
\usepackage{xcolor}

\definecolor{kbBack}{HTML}{F2F6FA}\definecolor{kbFrame}{HTML}{9DB4CC}\definecolor{kbTitle}{HTML}{52708F}
\definecolor{ibBack}{HTML}{FBF5E7}\definecolor{ibFrame}{HTML}{C9A961}\definecolor{ibTitle}{HTML}{7D5F1E}
\definecolor{wbBack}{HTML}{FCF2F0}\definecolor{wbFrame}{HTML}{D6A096}\definecolor{wbTitle}{HTML}{9E4F43}
\definecolor{dbBack}{HTML}{F4F7F3}\definecolor{dbFrame}{HTML}{AEC2AC}\definecolor{dbTitle}{HTML}{5F7F64}

\tcbset{softbox/.style={enhanced,breakable,boxrule=0.8pt,arc=2pt,left=3mm,right=3mm,top=2mm,bottom=2mm,coltitle=white,fonttitle=\bfseries,toptitle=0.6mm,bottomtitle=0.6mm}}
\newtcolorbox{knowledgebox}[1]{softbox,colback=kbBack,colframe=kbFrame,colbacktitle=kbTitle,title={#1}}
\newtcolorbox{importantbox}[1]{softbox,colback=ibBack,colframe=ibFrame,colbacktitle=ibTitle,title={#1}}
\newtcolorbox{warningbox}[1]{softbox,colback=wbBack,colframe=wbFrame,colbacktitle=wbTitle,title={#1}}
\newtcolorbox{dialoguebox}[1]{softbox,colback=dbBack,colframe=dbFrame,colbacktitle=dbTitle,title={#1}}

\lstset{language=Python,basicstyle=\ttfamily\small,keywordstyle=\color{blue!65!black},stringstyle=\color{red!60!black},commentstyle=\color{green!45!black},breaklines=true,frame=single,numbers=left,numberstyle=\tiny\color{gray},captionpos=b,extendedchars=true}
\hypersetup{colorlinks=true,linkcolor=blue!50!black,urlcolor=blue!60!black}
\setlength{\parindent}{2em}
\setlength{\parskip}{0.35em}
\setlength{\emergencystretch}{2em}
\renewcommand{\arraystretch}{1.25}

\newcommand{\makelecturetitle}{
\begin{titlepage}
\centering
\vspace*{0.8cm}
{\huge\bfseries \notetitle\par}
\vspace{0.6cm}
{\large \lecturers\par}
\vspace{0.25cm}
{\large \noteauthors\par}
\vspace{0.8cm}
\includegraphics[width=0.86\textwidth,height=0.44\textheight,keepaspectratio]{cover.jpg}
\vfill
\begin{tcolorbox}[width=0.92\textwidth,colback=black!2!white,colframe=black!55,sharp corners]
\textbf{Curso}: Stanford CS329A Self-Improving AI Agents (Autumn 2025)\par
\textbf{Fecha de la clase}: \lecturedate\quad
\textbf{Fecha de publicación del video}: 2026-08-03\par
\textbf{Fecha de elaboración de las notas}: \notedate\par
\textbf{Canal / duración}: Stanford Online\quad / \videoduration\par
\textbf{Enlace del video}: \href{\videourl}{\nolinkurl{\videourl}}\par
\textbf{Normas de elaboración}: \href{https://github.com/wdkns/wdkns-skills}{wdkns/wdkns-skills} (commit \texttt{39f1a04}) y las normas de calidad de este proyecto
\end{tcolorbox}
\end{titlepage}
\tableofcontents
\newpage
}

\newcommand{\videofigure}[4][0.95]{
\begin{figure}[H]
\centering
\includegraphics[width=#1\textwidth]{#2}
\caption{#3\protect\footnotemark}
\end{figure}
\footnotetext{Intervalo del video: #4.}
}
